\section{Second-order risk control}
The designated coefficient perturbations are unbiased. Together with bounded hidden features, this gives a small variance and a second-order mean bias even though the unconditioned prior has infinite coordinate variance.

\begin{lemma}[Pointwise perturbation moments]\label{lem:perturbation}
Let $g\in\calG_{m,B}$, and write $z_g$ for its raw linear readout. The comparison measure in~\eqref{eq:embedding-kl} satisfies
\begin{equation}
 \sup_x\E_\nu|z_w(x)-z_g(x)|^2
 \le S_{m,B}\epsilon^{2-\alpha},\qquad
 \sup_x\left|\E_\nu[z_w(x)-z_g(x)]\right|
 \le S_{m,B}\epsilon^{2-\alpha}.
 \label{eq:raw-moments}
\end{equation}
Here $S_{m,B}$ may be taken as in~\eqref{eq:resources}, independently of learner width.
\end{lemma}
\begin{proof}
First condition on a particular index assignment in the large-width construction; the small-width construction is already a product law. All selected incoming rows are independent of earlier layers. All hidden coordinates, selected or unselected, have magnitude at most one.

Put $\theta=\epsilon^{2-\alpha}$; then $\epsilon^2\le\theta$. Let $v_\ell$ be the largest second moment of a selected hidden neuron's error relative to its corresponding teacher value, and let $b_\ell$ be the largest absolute mean error. We fix an input throughout; all bounds are uniform over inputs and assignments.

In the first layer the preactivation error is centered and has second moment at most $(d+1)\epsilon^2/3$. At a later selected node, write the preactivation error as
\[
 \Delta=\sum_{j=1}^m a_j(h_j-h_j^*)+Z,
\]
where $a_j$ are the teacher coefficients, $h_j^*$ the teacher features, and $Z$ contains the designated coefficient perturbations and the background sum. Conditional on all preceding layers,
\begin{equation}
 \E[Z\mid\text{preceding layers}]=0,\qquad
 \E[Z^2\mid\text{preceding layers}]
 \le C_\alpha(m+1)\theta.
 \label{eq:centered-background}
\end{equation}
Indeed, the incoming coordinates are independent and centered after subtraction of their designated means. Their variances are bounded by Lemmas~\ref{lem:truncation} and~\ref{lem:active}, and all multiplying hidden features are bounded by one. This argument does not require the unselected hidden neurons to be close to zero.

The conditional centering removes the cross term, so
\[
 \E\Delta^2\le m^2B^2v_{\ell-1}+C_\alpha(m+1)\theta,
 \qquad |\E\Delta|\le mB b_{\ell-1}.
\]
Lipschitzness and Taylor's theorem therefore give
\begin{align}
 v_\ell&\le A^2\{m^2B^2v_{\ell-1}+C_\alpha(m+1)\theta\},\\
 b_\ell&\le AmB b_{\ell-1}
       +\frac J2\{m^2B^2v_{\ell-1}+C_\alpha(m+1)\theta\}.
 \label{eq:moment-recurrence}
\end{align}
The initial bounds are $v_1\le A^2(d+1)\theta/3$ and $b_1\le J(d+1)\theta/6$.

For completeness, set $a=AmB\ge1$ and $G=1+a$. Since $m+1\le2a$ and $G^2-a^2=1+2a$, induction in the first recurrence gives
\[
 v_\ell\le C_\alpha A^2(d+1)G^{2\ell-2}\theta.
\]
The second recurrence is absorbed by $G^2-a\ge a^2+a+1$, giving
\[
 b_\ell\le C_\alpha J(d+1)G^{2\ell-2}\theta.
\]
The same constants work at every depth: the gaps $G^2-a^2$ and $G^2-a$ absorb each additive term in the displayed recurrences, so no layer-dependent prefactor is introduced. The readout has the same conditional variance estimate, but no activation. Its second moment is at most $m^2B^2v_H+C_\alpha(m+1)\theta$, and its mean error is at most $mB b_H$. Both are bounded by a constant times $(1+A^2+J)(d+1)G^{2H}\theta$. This proves~\eqref{eq:raw-moments} for every assignment and hence also for their mixture.
\end{proof}

A Lipschitz estimate on the loss would only turn a mean-square error $\theta$ into a risk error of order $\sqrt\theta$. The next property preserves the second-order improvement, including at clipping thresholds.

\begin{lemma}[Squared loss after a bounded output link]\label{lem:semiconcavity}
For $|y|\le1$, let $\ell_y(z)=(\psi(z)-y)^2$. At every $z\in\R$ there is $d_y(z)$ such that, for every $u\in\R$,
\begin{equation}
 \ell_y(z+u)\le\ell_y(z)+d_y(z)u+C_\psi u^2,
 \qquad |d_y(z)|\le D_\psi,\qquad D_\psi+C_\psi\le B_\psi.
 \label{eq:semiconcavity}
\end{equation}
For clipping one can use $D_\psi=4$, $C_\psi=1$. For the smooth link class, one can use $D_\psi=4L_\psi$ and $C_\psi=L_\psi^2+2J_\psi$.
\end{lemma}
\begin{proof}
For clipping, $\ell_y(z)-z^2$ is concave: it is affine inside $[-1,1]$, has second derivative $-2$ outside, and its derivative jumps downward at both endpoints. Its supergradient inequality proves the claim, with $d_y(z)=2(z-y)$ for $|z|<1$ and zero for $|z|\ge1$.
For a smooth link, $|\ell_y'|\le4L_\psi$. Its derivative is Lipschitz with constant at most $2L_\psi^2+4J_\psi$, because $|\psi-y|\le2$. Taylor's inequality proves the displayed estimate.
\end{proof}

\begin{lemma}[A deterministic comparison for every sample]\label{lem:comparison}
For every $g\in\calG_{m,B}$, $0<r\le1$ and $N\ge m$, there is a measure $\nu_{g,r,N}\ll\mu_N$, independent of the training sample, satisfying
\begin{equation}
 \E_\nu(f_w(x)-y)^2\le(g(x)-y)^2+r
 \quad\text{for every }x\in[-1,1]^d,\ |y|\le1,
 \qquad \KL(\nu\|\mu_N)\le K_{m,B}r^{-q}.
 \label{eq:comparison}
\end{equation}
In particular the empirical-risk inequality holds for every sample with labels in $[-1,1]$, and the population-risk inequality holds for every law with labels in that interval. Section~\ref{sec:noisy-regression} uses it with the bounded reference label $f_0(x)$ to treat unbounded observed labels.
\end{lemma}
\begin{proof}
By Lemmas~\ref{lem:perturbation} and~\ref{lem:semiconcavity},
\[
 \E_\nu\ell_y(z_w(x))\le\ell_y(z_g(x))+B_\psi S_{m,B}\epsilon^{2-\alpha}.
\]
Choose $\epsilon=(r/(B_\psi S_{m,B}))^{1/(2-\alpha)}\le1$. Then $\epsilon^{-\alpha}=(B_\psi S_{m,B})^q r^{-q}$. Substitute this in~\eqref{eq:embedding-kl}. The additional logarithm of $1/r$ is bounded using $r^q\log(1/r)\le1/(eq)$, and the remaining terms are bounded by~\eqref{eq:complexity}.
\end{proof}

\begin{corollary}[Function-neighborhood mass]\label{cor:mass}
For any input law $D_X$, any $g\in\calG_{m,B}$, and $0<t\le1$, the same prior satisfies
\begin{equation}
 \mu_N\{\|f_w-g\|_{L^2(D_X)}\le t\}
 \ge \exp\{-C_\alpha K_{m,B}t^{-2q}\},
 \qquad N\ge m.
 \label{eq:function-mass}
\end{equation}
Thus the comparison argument also supplies an actual width-independent lower bound on a function neighborhood's prior probability.
\end{corollary}
\begin{proof}
Use Lemma~\ref{lem:comparison} with the noiseless law $Y=g(X)$ and $r=t^2/2$. Markov's inequality gives comparison probability at least $1/2$ to the displayed event. The binary relative-entropy inequality gives $\KL(\nu\|\mu_N)\ge\tfrac12\log(1/\mu_N(E))-\log2$, proving the claim after increasing the constant.
\end{proof}
